\documentclass[11pt]{article}

\usepackage[T1]{fontenc}
\usepackage{amsmath,amssymb,amsthm}
\usepackage[margin=1in]{geometry}
\usepackage[colorlinks=true,linkcolor=blue,urlcolor=blue]{hyperref}

\newtheorem{theorem}{Theorem}
\newtheorem{proposition}[theorem]{Proposition}
\newtheorem{corollary}[theorem]{Corollary}
\newtheorem{conjecture}[theorem]{Conjecture}
\theoremstyle{remark}
\newtheorem{remark}[theorem]{Remark}

\newcommand{\R}{\mathbb R}
\newcommand{\Pp}{\mathbb P}
\newcommand{\norm}[1]{\left\|#1\right\|}
\newcommand{\abs}[1]{\left|#1\right|}
\newcommand{\ip}[2]{\left\langle #1,#2\right\rangle}
\newcommand{\supp}{\operatorname{supp}}
\newcommand{\Var}{\operatorname{Var}}
\setlength{\emergencystretch}{3em}

\title{A Finite Weighted Transfer Budget Forces\\
q4 Spectral-Shell Broadening}
\author{John Lawson and Codex}
\date{25 July 2026}

\begin{document}
\maketitle

\begin{abstract}
This report advances the conditional, energy-efficient q4 Type-II
entrance without resolving a Clay alternative.  The radial nonlinear
transfer has a finite same-trajectory norm:
\[
 \int_0^{s_0}\int_0^\infty
 \lambda^{-1/2}\,d|\mu_s|(\lambda)\,ds<\infty.
\]
It follows from \(\dot H^{-1}\)--\(L^2\) duality and the finite
\(L_s^2\) ledgers for weak-\(L^3\) amplitude and enstrophy amplitude.
Differentiating a fixed positive energy quantile gives an exact
spectral-speed law.  On a descending q4 shell it forces the effective
local width \(W_\eta\) to satisfy
\[
 \int_0\frac{\sqrt{L(s)}}{sW_\eta(s)}\,ds<\infty.
\]
Thus persistent width \(W_\eta\lesssim\sqrt L\), including any
nondegenerate smoothing of adjacent Pell shells, is impossible.

This is a genuine conditional reduction, but it does not close q4.
We construct a smooth positive spectral density of fixed relative width
\(W_\eta\asymp L\).  Its characteristic flux is strictly positive at
every positive heat age for all sufficiently small times.  It obeys the
exact energy law, matches every current q4 scalar and Gaussian budget,
has asymptotically complete early-history reuse, and pays the new
weighted transfer charge finitely.  The remaining known unknown is
therefore a spatial, angular, phase, or genealogical cost for broad
same-trajectory transfer.  The countermodel is a spectrum rather than a
velocity field.  All novel deductions await external review.
\end{abstract}

\section{Conditional entrance and spectral balance}

Let \(\nu>0\).  On \(0<s<s_0\), suppose that \(v\) and \(U\) are smooth,
divergence-free fields on \(\R^3\) satisfying
\begin{equation}
\label{eq:evolution}
 \partial_s v-\nu\Delta v+F=0,
 \qquad
 F:=\Pp((U\cdot\nabla)v).
\end{equation}
The conditional q4 entrance established in the preceding rounds supplies
\begin{equation}
\label{eq:energy}
 \norm{v(s)}_2^2+
 2\nu\int_0^s\norm{\nabla v(a)}_2^2\,da=d>0,
\end{equation}
\begin{equation}
\label{eq:heat-boundary}
 e^{\nu r\Delta}v(s)\longrightarrow0
 \quad\hbox{strongly in \(L^2\) for every fixed \(r>0\),}
\end{equation}
as \(s\downarrow0\).  Write
\begin{equation}
\label{eq:MY}
 M(s):=\norm{U(s)}_{L^{3,\infty}},
 \qquad
 Y(s):=\norm{\nabla v(s)}_2,
\end{equation}
and assume the proved entrance ledgers
\begin{equation}
\label{eq:L2-ledgers}
 M,Y\in L^2(0,s_0).
\end{equation}

Push the Fourier energy and nonlinear pairing forward by
\(\lambda=2\nu|\xi|^2\):
\begin{align}
\label{eq:epsilon}
 \varepsilon_s(B)
 &:=
 \frac12\int_{\{2\nu|\xi|^2\in B\}}
 |\widehat v(s,\xi)|^2\,d\xi,\\
\label{eq:mu}
 \mu_s(B)
 &:=
 -\operatorname{Re}
 \int_{\{2\nu|\xi|^2\in B\}}
 \widehat F(s,\xi)\cdot\overline{\widehat v(s,\xi)}\,d\xi.
\end{align}
The previous round proved
\begin{equation}
\label{eq:characteristic}
 \mathcal E(r,s):=
 \frac12\norm{e^{\nu r\Delta}v(s)}_2^2
 =\int_0^\infty e^{-r\lambda}\,d\varepsilon_s(\lambda),
\end{equation}
\begin{equation}
\label{eq:flux}
 \Pi_r(s):=\int_0^\infty e^{-r\lambda}\,d\mu_s(\lambda)
 =(\partial_s-\partial_r)\mathcal E(r,s),
\end{equation}
and, distributionally,
\begin{equation}
\label{eq:measure-balance}
 \mu_s=\partial_s\varepsilon_s+\lambda\varepsilon_s.
\end{equation}
When the radial measures have densities, write
\begin{equation}
\label{eq:density-balance}
 \varepsilon_s(d\lambda)=e(\lambda,s)\,d\lambda,
 \qquad
 \mu_s(d\lambda)=m(\lambda,s)\,d\lambda,
 \qquad
 \partial_se=m-\lambda e.
\end{equation}
Only the quantile statements below require the density and
nondegeneracy hypotheses.

\section{The finite weighted transfer theorem}

Define the weighted radial transfer norm
\begin{equation}
\label{eq:q-def}
 \mathfrak q(s):=
 \int_0^\infty\lambda^{-1/2}\,d|\mu_s|(\lambda).
\end{equation}

\begin{theorem}[Finite weighted nonlinear transfer]
\label{thm:weighted}
For almost every \(s\in(0,s_0)\),
\begin{equation}
\label{eq:q-bound}
 \mathfrak q(s)
 \le
 (2\nu)^{-1/2}
 \norm{F(s)}_{\dot H^{-1}}\norm{v(s)}_2
 \lesssim
 \nu^{-1/2}\sqrt d\,M(s)Y(s).
\end{equation}
Consequently,
\begin{equation}
\label{eq:q-L1}
 \int_0^{s_0}\mathfrak q(s)\,ds<\infty.
\end{equation}
\end{theorem}

\begin{proof}
Total variation cannot increase under a measurable pushforward.  Hence
\begin{align}
 \mathfrak q(s)
 &\le
 \int_{\R^3}
 (2\nu|\xi|^2)^{-1/2}
 |\widehat F(s,\xi)|\,|\widehat v(s,\xi)|\,d\xi \notag\\
 &\le
 (2\nu)^{-1/2}
 \norm{F(s)}_{\dot H^{-1}}\norm{v(s)}_2.
\label{eq:pushforward-CS}
\end{align}
For a solenoidal test field \(\psi\), integration by parts, Lorentz
Hölder, and the Lorentz refinement of Sobolev give
\begin{align}
 \abs{\ip{F}{\psi}}
 &=
 \abs{\int_{\R^3}U_jv_i\,\partial_j\psi_i\,dx}\notag\\
 &\lesssim
 \norm U_{L^{3,\infty}}
 \norm v_{L^{6,2}}
 \norm{\nabla\psi}_2
 \lesssim
 MY\norm{\nabla\psi}_2.
\label{eq:Hminus}
\end{align}
Because \(F\) is solenoidal, the same estimate for arbitrary tests
follows by applying the bounded Leray projector.  Thus
\(\norm F_{\dot H^{-1}}\lesssim MY\).  Equation
\eqref{eq:energy} gives \(\norm v_2\le\sqrt d\), proving
\eqref{eq:q-bound}.  Finally,
\[
 \int_0^{s_0}M(s)Y(s)\,ds
 \le
 \norm M_{L^2(0,s_0)}\norm Y_{L^2(0,s_0)}<\infty,
\]
which proves \eqref{eq:q-L1}.
\end{proof}

\begin{remark}
The unweighted radial variation costs the already divergent q4 density
\(MY^2\).  The factor \(\lambda^{-1/2}\) is one heat length; it transfers
one derivative to the finite kinetic-energy factor and lowers the cost
to the integrable product \(MY\).
\end{remark}

\section{Exact spectral-quantile speed}

Fix \(0<\eta<d/2\).  Whenever it exists, define the low-frequency
energy quantile \(\Lambda_\eta(s)\) by
\begin{equation}
\label{eq:quantile}
 \int_0^{\Lambda_\eta(s)}e(\lambda,s)\,d\lambda=\eta.
\end{equation}
The heat boundary ensures that every such fixed quantile escapes to
infinity.  Indeed, for fixed \(r,\Lambda>0\),
\begin{equation}
\label{eq:low-frequency}
 \int_0^\Lambda e(\lambda,s)\,d\lambda
 \le e^{r\Lambda}\mathcal E(r,s)\longrightarrow0.
\end{equation}
Since the total energy tends to \(d/2\), \(\Lambda_\eta(s)\to\infty\).

\begin{theorem}[Fixed-energy quantile law]
\label{thm:quantile}
Suppose that \(e\) and \(m\) are locally integrable, that \(e\) is
continuous and positive at \(\lambda=\Lambda_\eta(s)\), and that the
quantile is differentiable.  Then
\begin{equation}
\label{eq:quantile-speed}
 e(\Lambda_\eta,s)\Lambda_\eta'
 =
 \int_0^{\Lambda_\eta}\lambda e(\lambda,s)\,d\lambda
 -
 \int_0^{\Lambda_\eta}m(\lambda,s)\,d\lambda.
\end{equation}
On
\(\mathcal D_\eta:=\{s:\Lambda_\eta'(s)<0\}\),
\begin{equation}
\label{eq:quantile-charge}
 \int_{\mathcal D_\eta}
 e(\Lambda_\eta(s),s)
 \frac{|\Lambda_\eta'(s)|}{\sqrt{\Lambda_\eta(s)}}\,ds
 \le
 \int_0^{s_0}\mathfrak q(s)\,ds<\infty.
\end{equation}
\end{theorem}

\begin{proof}
Differentiate \eqref{eq:quantile} and use
\eqref{eq:density-balance}:
\[
 0=
 \int_0^{\Lambda_\eta}(m-\lambda e)\,d\lambda
 +e(\Lambda_\eta,s)\Lambda_\eta'.
\]
This is \eqref{eq:quantile-speed}.  If
\(\Lambda_\eta'<0\), then
\begin{align*}
 e(\Lambda_\eta,s)|\Lambda_\eta'|
 &=
 \int_0^{\Lambda_\eta}m\,d\lambda
 -\int_0^{\Lambda_\eta}\lambda e\,d\lambda\\
 &\le
 \int_0^{\Lambda_\eta}|m|\,d\lambda
 \le
 \sqrt{\Lambda_\eta}\,\mathfrak q(s).
\end{align*}
Division by \(\sqrt{\Lambda_\eta}\), integration, and
Theorem~\ref{thm:weighted} prove \eqref{eq:quantile-charge}.
\end{proof}

\subsection*{A smooth Gaussian quantile}

There is also a density-free heat formulation.  For \(0<\theta<1\),
let \(r_\theta(s)\) solve
\begin{equation}
\label{eq:heat-quantile}
 \mathcal E(r_\theta(s),s)=\theta\mathcal E(0,s)
\end{equation}
and put
\begin{equation}
\label{eq:D}
 \mathcal D(r,s):=-\partial_r\mathcal E(r,s)
 =\nu\norm{\nabla e^{\nu r\Delta}v(s)}_2^2.
\end{equation}
Since \(\Pi_0=0\), differentiating \eqref{eq:heat-quantile} and using
\eqref{eq:flux} gives the exact identity
\begin{equation}
\label{eq:heat-speed}
 \boxed{
 r_\theta'(s)
 =
 \frac{\Pi_{r_\theta(s)}(s)+\theta\mathcal D(0,s)}
 {\mathcal D(r_\theta(s),s)}
 -1.}
\end{equation}
It separates nonlinear shell motion from preferential viscous loss.

\section{The q4 narrowing obstruction}

Let
\begin{equation}
\label{eq:L}
 \ell(s):=\log(e/s),
 \qquad
 L(s):=s^{-9/11}\ell(s)^{2/11},
\end{equation}
and define the effective local quantile width
\begin{equation}
\label{eq:W}
 W_\eta(s):=\frac{\eta}{e(\Lambda_\eta(s),s)}.
\end{equation}

\begin{corollary}[Persistent narrow shells are impossible]
\label{cor:narrow}
Assume, for all sufficiently small \(s\), that
\begin{equation}
\label{eq:q4-descent}
 \Lambda_\eta(s)\asymp L(s),
 \qquad
 -\Lambda_\eta'(s)\gtrsim\frac{L(s)}s.
\end{equation}
Then
\begin{equation}
\label{eq:width-charge}
 \int_0
 \frac{\sqrt{L(s)}}{sW_\eta(s)}\,ds<\infty.
\end{equation}
In particular, an eventual bound
\begin{equation}
\label{eq:narrow}
 W_\eta(s)\lesssim\sqrt{L(s)}
\end{equation}
is impossible.
\end{corollary}

\begin{proof}
On the descending set, the integrand in
\eqref{eq:quantile-charge} satisfies
\[
 e(\Lambda_\eta,s)
 \frac{|\Lambda_\eta'|}{\sqrt{\Lambda_\eta}}
 \gtrsim
 \frac{\eta\sqrt{L(s)}}{sW_\eta(s)}.
\]
This proves \eqref{eq:width-charge}.  Under \eqref{eq:narrow}, its
integrand is bounded below by a positive multiple of \(1/s\), a
contradiction.
\end{proof}

\begin{remark}[Disposition of the Pell snapshots]
The exact Pell family transfers between
\[
 \lambda_m=8\nu m^2,
 \qquad
 \lambda_m+2\nu.
\]
Any nondegenerate \(O(1)\)-width smoothing that carries a fixed positive
energy quantile has \(W_\eta=O(1)\ll\sqrt{\lambda_m}\), and so cannot be
concatenated continuously along \eqref{eq:q4-descent}.  The exact torus
snapshots are atomic and do not satisfy the density hypothesis in
Theorem~\ref{thm:quantile}; their instantaneous transfer theorem remains
valid.  What is closed is their identification with one persistently
narrow moving quantile.
\end{remark}

More generally, the power ansatz \(W_\eta=L^\gamma\) requires
\(\gamma>1/2\) at power level.  The endpoint \(\gamma=1/2\) already
incurs the logarithmic divergence \(\int ds/s\).

\section{A sharp broad-shell survivor}

Choose \(0<a<1<b\) and
\begin{equation}
\label{eq:phi}
 \phi\in C_c^\infty((0,\infty)),
 \quad
 \supp\phi=[a,b],
 \quad
 \phi>0\ \hbox{on }(a,b),
 \quad
 \int\phi=\int y\phi(y)\,dy=1.
\end{equation}
Define
\begin{equation}
\label{eq:A}
 A(s):=\frac d2
 \exp\left(-\int_0^sL(\sigma)\,d\sigma\right),
 \qquad A'=-AL,
\end{equation}
\begin{equation}
\label{eq:ephi}
 e_\phi(\lambda,s):=
 \frac{A(s)}{L(s)}
 \phi\left(\frac{\lambda}{L(s)}\right).
\end{equation}
Let
\begin{equation}
\label{eq:Phi}
 \Phi(x):=\int_0^\infty e^{-xy}\phi(y)\,dy,
 \qquad
 G(x):=-\Phi'(x)>0,
 \qquad
 \alpha(s):=-\frac{sL'(s)}{L(s)}.
\end{equation}

\begin{proposition}[Exact broad-shell characteristic]
\label{prop:broad}
The density \eqref{eq:ephi} is positive, has mass \(A\) and mean
frequency \(L\):
\begin{equation}
\label{eq:mass-mean}
 \int e_\phi\,d\lambda=A,
 \qquad
 \int\lambda e_\phi\,d\lambda=AL.
\end{equation}
Its heat energy and characteristic flux are
\begin{equation}
\label{eq:Ephi}
 \mathcal E_\phi(r,s)=A(s)\Phi(rL(s)),
\end{equation}
\begin{equation}
\label{eq:Piphi}
 \boxed{
 \Pi^\phi_r(s)
 =
 \frac{A(s)}s
 \left[
 \alpha(s)xG(x)
 -sL(s)\bigl(\Phi(x)+\Phi'(x)\bigr)
 \right],\qquad x=rL(s).}
\end{equation}
The corresponding transfer density is
\begin{equation}
\label{eq:mphi}
 \boxed{
 m_\phi(\lambda,s)
 =
 A(s)\left[
 (y-1)\phi(y)
 -\frac{L'(s)}{L(s)^2}
 \bigl(\phi(y)+y\phi'(y)\bigr)
 \right],\qquad y=\frac{\lambda}{L(s)}.}
\end{equation}
For every sufficiently small \(s>0\),
\begin{equation}
\label{eq:positive-flux}
 \Pi^\phi_r(s)>0
 \qquad\hbox{for every \(r>0\).}
\end{equation}
\end{proposition}

\begin{proof}
The change of variables \(\lambda=Ly\), together with
\eqref{eq:phi}, proves \eqref{eq:mass-mean}.  It also gives
\eqref{eq:Ephi}.  Differentiating,
\[
 \partial_s\mathcal E_\phi=A'\Phi+ArL'\Phi',
 \qquad
 \partial_r\mathcal E_\phi=AL\Phi'.
\]
Use \(A'=-AL\), \(x=rL\), \(\Phi'=-G\), and the definition of
\(\alpha\) to obtain
\((\partial_s-\partial_r)\mathcal E_\phi=\Pi^\phi\) in the form
\eqref{eq:Piphi}.  Finally, differentiating \eqref{eq:ephi} and using
\(m_\phi=\partial_se_\phi+\lambda e_\phi\) gives
\eqref{eq:mphi}.

For \eqref{eq:positive-flux}, introduce the exponentially tilted mean
\begin{equation}
\label{eq:tilted-mean}
 \overline y(x):=\frac{G(x)}{\Phi(x)}.
\end{equation}
With \(\operatorname{Var}_x\) denoting variance under the probability
density \(e^{-xy}\phi(y)/\Phi(x)\), it satisfies
\begin{equation}
\label{eq:tilted-variance}
 \overline y(0)=1,
 \qquad
 \overline y'(x)=-\operatorname{Var}_x(y)<0.
\end{equation}
Consequently,
\begin{equation}
\label{eq:H-positive}
 H(x):=\Phi(x)+\Phi'(x)
 =\Phi(x)(1-\overline y(x))>0
 \qquad(x>0).
\end{equation}
The quotient \(H(x)/(xG(x))\) tends to the variance of \(\phi\) as
\(x\downarrow0\), and it tends to zero as \(x\to\infty\), since the
tilted mean tends to \(a>0\).  Thus
\begin{equation}
\label{eq:Cphi}
 C_\phi:=\sup_{x>0}\frac{H(x)}{xG(x)}<\infty.
\end{equation}
Because \(\alpha(s)\to9/11\) and \(sL(s)\to0\),
\eqref{eq:Piphi} gives
\[
 \Pi^\phi_r(s)
 \ge
 \frac{A(s)}s xG(x)
 \bigl(\alpha(s)-sL(s)C_\phi\bigr)>0
\]
for every \(x>0\) and all sufficiently small \(s\).
\end{proof}

\begin{proposition}[The new budget is sharp against broad shells]
\label{prop:sharp}
The broad shell satisfies
\begin{equation}
\label{eq:qphi}
 \mathfrak q_\phi(s)
 \lesssim
 A(s)\left(
 \sqrt{L(s)}
 +\frac{|L'(s)|}{L(s)^{3/2}}
 \right)
 \asymp
 A(s)\left(
 \sqrt{L(s)}
 +\frac1{s\sqrt{L(s)}}
 \right).
\end{equation}
Both terms are integrable at zero.  Moreover, every fixed
\(0<\eta<d/2\) has, for small \(s\),
\begin{equation}
\label{eq:broad-quantile}
 \Lambda_\eta(s)=L(s)z_\eta(s),
 \qquad
 W_\eta(s)\asymp L(s),
 \qquad
 \Lambda_\eta'(s)=z_\eta(s)L'(s)(1+o(1)),
\end{equation}
and its charge is of order
\begin{equation}
\label{eq:broad-charge}
 \frac1{s\sqrt{L(s)}}\in L^1(0,s_0).
\end{equation}
\end{proposition}

\begin{proof}
Substitute \(y=\lambda/L\) in \eqref{eq:mphi}.  Compact support away
from zero gives
\[
 \int_0^\infty\lambda^{-1/2}|m_\phi|\,d\lambda
 \lesssim
 A\left(\sqrt L+\frac{|L'|}{L^{3/2}}\right).
\]
Since \(|L'|/L\asymp1/s\), the two singular factors are
\[
 s^{-9/22}\ell(s)^{1/11},
 \qquad
 s^{-13/22}\ell(s)^{-1/11},
\]
and both are integrable.

Let \(z_\eta(s)\in(a,b)\) be determined by
\[
 \int_a^{z_\eta(s)}\phi(y)\,dy=\frac{\eta}{A(s)}.
\]
It converges to an interior point because \(A(s)\to d/2\) and
\(0<\eta<d/2\).  Thus \(\phi(z_\eta)\) stays positive and
\[
 \Lambda_\eta=Lz_\eta,
 \qquad
 e_\phi(\Lambda_\eta,s)=\frac A L\phi(z_\eta),
\]
which proves \(W_\eta\asymp L\).  Differentiating the defining equation
gives \(z_\eta'=O(L)\).  Because \(L^2/|L'|\asymp sL\to0\),
\[
 \Lambda_\eta'=z_\eta L'+Lz_\eta'
 =z_\eta L'(1+o(1)).
\]
Substitution in the quantile charge yields
\eqref{eq:broad-charge}.
\end{proof}

\begin{proposition}[All previous q4 scalar tests survive]
\label{prop:q4}
Assign the broad spectrum the scalar ledgers
\begin{equation}
\label{eq:ledger-def}
 M_\phi(s):=\frac1{sL(s)},
 \qquad
 Y_\phi(s)^2:=\frac{A(s)L(s)}{\nu}.
\end{equation}
Then
\begin{align}
 \int_0^hM_\phi(s)^2\,ds
 &\asymp h^{7/11}\ell(h)^{-4/11},\label{eq:Mtail}\\
 \int_0^hY_\phi(s)^2\,ds
 &\asymp h^{2/11}\ell(h)^{2/11},\label{eq:Ytail}\\
 M_\phi(s)Y_\phi(s)^2
 &=\frac{A(s)}{\nu s}.\label{eq:MY2}
\end{align}
The flux obeys
\begin{equation}
\label{eq:pointwise}
 |\Pi^\phi_r(s)|
 \lesssim r^{-1/2}M_\phi(s)Y_\phi(s),
\end{equation}
\begin{equation}
\label{eq:variation}
 \Var_{r>0}\Pi^\phi_r(s)\lesssim\frac{A(s)}s.
\end{equation}
For every fixed \(0<q<1\),
\begin{equation}
\label{eq:early}
 \frac{\displaystyle
 \int_0^{q\tau}\Pi^\phi_{\tau-s}(s)\,ds}
 {\displaystyle
 \int_0^\tau\Pi^\phi_{\tau-s}(s)\,ds}
 \longrightarrow1
 \qquad(\tau\downarrow0).
\end{equation}
\end{proposition}

\begin{proof}
Equations \eqref{eq:Mtail}--\eqref{eq:MY2} are the elementary
power--log integrals associated with \eqref{eq:L}.  For the heat bounds,
put \(x=rL(s)\).  The profiles
\[
 xG(x),
 \qquad
 \Phi(x)+\Phi'(x)
\]
are smooth and exponentially decreasing.  The second vanishes at zero
because \(\int\phi=\int y\phi=1\).  Therefore
\[
 x^{1/2}|xG(x)|
 +x^{1/2}|\Phi(x)+\Phi'(x)|\lesssim1.
\]
The motion coefficient in \eqref{eq:Piphi} is \(A/s\); the dispersion
coefficient is \(AL\), smaller by \(sL\to0\).  Comparison with
\[
 r^{-1/2}M_\phi Y_\phi
 =
 \sqrt{\frac A\nu}\,\frac{x^{-1/2}}s
\]
proves \eqref{eq:pointwise}, with constants allowed to depend on the
fixed entrance parameters.  Finite variation of the two profiles proves
\eqref{eq:variation}.

The characteristic identity and the fixed heat boundary give
\[
 \int_0^\tau\Pi^\phi_{\tau-s}(s)\,ds=A(\tau).
\]
On \(q\tau\le s\le\tau\), the integrated motion term is
\(O_q(\tau L(q\tau))\).  Since the dispersion profile vanishes linearly
at zero, its integral is
\(O_q((\tau L(q\tau))^2)\).  Both vanish as \(\tau\downarrow0\), while
\(A(\tau)\to d/2\), proving \eqref{eq:early}.
\end{proof}

\section{Novel findings, closures, and open obligations}

\subsection*{Robust conditional deductions, pending external review}

\begin{enumerate}
 \item The actual radial nonlinear transfer has the finite
 \(\lambda^{-1/2}\)-weighted total-variation budget
 \eqref{eq:q-L1}.
 \item A descending fixed-energy spectral quantile obeys the exact
 speed law \eqref{eq:quantile-speed} and finite charge
 \eqref{eq:quantile-charge}.
 \item On the q4 schedule, persistent local width
 \(W_\eta\lesssim\sqrt L\) is impossible.
 \item This closes a nondegenerate dynamic concatenation of fixed-energy
 adjacent Pell shells while preserving the exact atomic snapshots.
 \item A smooth fixed-relative-width positive spectrum has strictly
 positive heat flux at every positive heat age for all sufficiently
 small times, satisfies every current radial, temporal, and scalar q4
 test, and pays the new charge finitely.  Thus the width exponent
 \(1/2\) is sharp against the present information.
\end{enumerate}

\subsection*{Closed shortcuts}

\begin{enumerate}
 \item The ideal delta-shell survivor is too narrow as stated.
 \item Adjacent-shell Pell snapshots cannot simply be interpreted as one
 continuously descending fixed-energy q4 quantile.
 \item The finite weighted-transfer theorem alone cannot close q4,
 because a broad relative shell survives it.
 \item Positivity, radial moments, Gaussian variation, and early-history
 reuse remain insufficient without a same-trajectory broad-shell cost.
\end{enumerate}

\subsection*{Things left to prove}

\begin{enumerate}
 \item Resolve the quadratic convolution into angular caps and derive a
 finite mode-count, phase, or parent-frequency charge for
 \(W_\eta\gg\sqrt L\).
 \item Test that charge against smooth fixed-relative-width dilations of
 the exact Pell family.
 \item Couple radial broadening to spatial concentration and the
 Hausdorff-one-null terminal energy measure.
 \item Use the strong-trace drift to constrain coherent steering of
 broad triadic transfer.
 \item Treat intermittent narrowing in place of an eventual pointwise
 width hypothesis.
 \item Treat slower clocks, divergent normalised energy, R3B, and the
 other Clay alternatives separately.
\end{enumerate}

\begin{conjecture}[No dynamically broad q4 descent]
No same-trajectory full-defect q4 entrance with finite unweighted
dissipation can keep a fixed positive energy quantile descending at
\[
 \Lambda_\eta(s)\asymp
 s^{-9/11}\ell(s)^{2/11}
\]
with the broad width \(W_\eta(s)\gg\sqrt{\Lambda_\eta(s)}\) required by
Theorem~\ref{thm:quantile}, while regenerating the full boundary defect
through positive early-history triadic work.
\end{conjecture}

\begin{remark}[Epistemic boundary]
The conjecture is open.  The broad shell is a positive spectral
countermodel, not a velocity field, and a dilated broad triad is not one
Navier--Stokes evolution.  This report proves no q4 closure, energy
equality, regularity theorem, breakdown theorem, or Clay alternative.
\end{remark}

\end{document}
